\documentclass[11pt]{article}
\usepackage{amsmath,amssymb,amsfonts}
\usepackage{geometry}
\usepackage{enumitem}
\usepackage{color}
\usepackage{hyperref}
\usepackage{booktabs}

\geometry{letterpaper, margin=1in}

\definecolor{addressed}{rgb}{0,0.5,0}
\definecolor{partial}{rgb}{0.8,0.4,0}
\definecolor{pending}{rgb}{0.8,0,0}

\title{\textbf{Response to Reviewers --- Round 2}\\[0.5em]
\large Manuscript: Flux-Preserving Adaptive Finite State Projection\\
for Multiscale Stochastic Reaction Networks\\[0.5em]
\normalsize Submitted to: SIAM Multiscale Modeling and Simulation}

\author{Aditya Dendukuri, Shivkumar Chandrasekaran, and Linda Petzold\\
Department of Computer Science\\
University of California, Santa Barbara}

\date{\today}

\begin{document}

\maketitle

\noindent Dear Editor and Reviewers,

\medskip

Thank you for the continued careful evaluation of our manuscript. We are
grateful for the constructive feedback in this second round, and we have
addressed each point raised by the referees. Below we respond point by point.

%======================================================================
\section*{Response to Referee 1}
%======================================================================

\noindent\textbf{Comment 1.1:} \textit{The ``total flux'' of a reduced state
space is defined as the sum of the flux of each state in the subset. The
majority of this total flux would then be the propensities to move between two
states already contained in the reduced state space. If the goal of the
algorithm is to avoid loss of states that move to other metastable regions,
would it make more sense to limit the flux to propensities of moving outside
the reduced state space?}

\textbf{Response:}
% TODO: Draft response about boundary-only flux vs total flux in pruning

\medskip

\noindent\textbf{Comment 1.2 (minor):} \textit{Line 140: the term for flux
$w(x)$ is used before being defined.}

\textbf{Response:} We have added an inline definition of $w(x)$ at its first
use in Remark~2.3, so that the reader does not need to look ahead to
Section~3 for the formal definition.

\medskip

\noindent\textbf{Comment 1.3 (minor):} \textit{Fig.\ 2: As the flux is central
to your paper, it would be helpful to visualize it for the readers. A natural
place would be in Fig.\ 2A, where the fluxes can be illustrated on the arrows.}

\textbf{Response:}
% TODO: Add flux labels to Fig 2 TikZ diagram and describe here

%======================================================================
\section*{Response to Referee 2}
%======================================================================

\subsection*{Theoretical Issues}

\noindent\textbf{Comment 2.1:} \textit{There is a significant notational error
in Eq.\ (4.4) regarding the definition of boundary fluxes. Given the stated
column convention, $\Phi_{\mathrm{out}}$ should involve $\mathbf{A}_{J'J}$
rather than $\mathbf{A}_{JJ'}$.}

\textbf{Response:} We thank the referee for catching this. We have added an
explicit clarification of the column convention immediately after Eq.~(4.4):
the block $\mathbf{A}_{J'J}$ has rows indexed by $J'$ and columns indexed
by~$J$, so that the product $\mathbf{A}_{J'J}\boldsymbol{p}_J$ gives the rate
at which probability flows \emph{out} of~$J$ into~$J'$. This convention is now
stated in the text and the notation is consistent throughout Section~4.

\medskip

\noindent\textbf{Comment 2.2:} \textit{The derivation for the adaptive
time-step rule ($\Delta t = \varepsilon_{\Delta t}/\Phi_{\max}$) relies on the
inequality $\Phi_{\mathrm{in}} + \Phi_{\mathrm{out}} \leq \Phi_{\max}$, for
which I fail to see the justification without a multiplicative constant that
accounts for the boundary's size or structure.}

\textbf{Response:} The referee is correct that the original inequality
$\Phi_{\mathrm{in}} + \Phi_{\mathrm{out}} \leq \Phi_{\max}$ was not justified.
We have replaced it with the trivially valid bound
$\Phi_{\mathrm{out}} \leq \Phi_{\mathrm{total}}$, which holds because the
boundary outflux is a subset of the total system flux. The adaptive time-step
rule in Corollary~4.7 now reads
$\Delta t \leq \varepsilon_{\Delta t}/\Phi_{\mathrm{total}}$, and the proof
follows immediately from this containment. This correction propagates cleanly
through the subsequent corollaries without affecting the structure of the
error bounds.

\medskip

\noindent\textbf{Comment 2.3:} \textit{The algorithm (3.1 and 3.3) performs
renormalization after pruning, yet the theoretical analysis in Section~4 is
built upon a compressed (reflecting) generator. The relationship between these
two treatments needs explicit clarification to ensure the error bounds remain
valid for the implemented code.}

\textbf{Response:} We appreciate this question, as it highlights a point that
deserved clarification. Upon reflection, there is in fact \emph{no gap} between
the theory and the implementation, because the implementation directly
constructs the compressed generator---not the truncated generator followed by
renormalization.

Specifically, Algorithm~3.2 (forward-enumeration matrix construction) builds
the generator column by column. For each source state $x \in J$ and each
reaction~$k$, the propensity $\alpha_k(x)$ contributes to the off-diagonal
entry at row $y = x + \boldsymbol{\nu}_k$ \emph{only if} $y \in J$. The
diagonal entry for column~$x$ is then set to the negative of the sum of these
included off-diagonal entries. Transitions whose destination falls outside~$J$
are excluded from both the off-diagonal and the diagonal. Consequently, the
column sums are exactly zero by construction. This is precisely the defining
property of the compressed generator $\widetilde{\mathbf{A}}_{JJ}$
(Proposition~4.1, property~1), and the matrix exponential
$e^{\widetilde{\mathbf{A}}_{JJ}\Delta t}$ therefore conserves
$\boldsymbol{1}^\top \boldsymbol{p}$ exactly.

The renormalization $\boldsymbol{p} \gets \boldsymbol{p}/\|\boldsymbol{p}\|_1$
in Algorithm~3.1 (Stage~3) serves a \emph{different} purpose: it compensates
for the probability mass discarded when low-probability states are removed
during pruning, not for any leakage from the generator. The pruned mass is at
most~$\alpha$ by the quantile selection in Stage~1.

We have added Remark~4.4 to the revised manuscript to make this distinction
explicit, stating that the implementation uses the compressed generator and that
the renormalization addresses pruning losses only. The error analysis in
Section~4, which is formulated entirely in terms of
$\widetilde{\mathbf{A}}_{JJ}$, therefore applies directly to the implemented
algorithm without any additional approximation gap.

\medskip

\noindent\textbf{Comment 2.4:} \textit{Estimating $w_{\max}$: Corollary~4.7
uses $w_{\max}$ (the maximum exit rate of pruned states), but it is unclear how
this value is robustly estimated on a truncated domain since it involves states
outside the current active set.}

\textbf{Response:} We have clarified this in the revised manuscript. The key
observation is that the exit rate $w(x) = \sum_k \alpha_k(x)$ is computed for
every state $x$ in the active set \emph{before} pruning takes place (it is
needed for the flux computation in Algorithm~3.1, Stage~2). Since pruning only
removes a subset of the current active states, $w_{\max}$ over the pruned
states is available as a byproduct of the flux computation, with no need to
evaluate $w(x)$ at states outside the active set. Corollary~4.7 now states this
explicitly.

\subsection*{Empirical Issues}

\noindent\textbf{Comment 2.5:} \textit{The manuscript lacks head-to-head
comparisons against strong adaptive FSP baselines, such as ETFSP or other
time-varying variants.}

\textbf{Response:} We have added Section~5.6 to the revised manuscript, which
provides a head-to-head comparison of the flux-based adaptive FSP against the
Krylov-FSP-SSA method of Vo and Sidje~[Sidje2015], reviewed by Dinh and
Sidje~[Dinh2016]. This baseline shares the same expand--evolve--prune skeleton
and Krylov matrix exponential integrator, differing in pruning (probability
threshold without flux protection), time stepping (mass-conservation
accept/reject rather than flux-adaptive), and expansion (SSA-driven with
reachability enumeration). The comparison isolates the effect of our two main
contributions---flux-aware pruning and flux-adaptive time stepping---on three
benchmark systems.

On the Oregonator, the flux-adaptive time stepping provides a decisive
advantage: the adaptive FSP completes the simulation in 41\,s using 28
snapshots with $|S| \leq 1{,}058$, while the Krylov-FSP-SSA method reaches only
9\% of the time horizon after 10{,}000 iterations and 438\,s, with the state
space growing to 15{,}850 states. On the bottleneck system, the flux-based
pruning preserves the essential $B = 1$ bottleneck state (which carries
probability $\sim 10^{-6}$, below the Krylov-FSP-SSA threshold) and correctly
tracks $\langle C \rangle = 357$, while the baseline produces
$\langle C \rangle = 18.7$ after severing the bottleneck pathway. On the
Robertson system over $[0, 10]$, both methods complete, with the Krylov-FSP-SSA
method faster due to the matrix exponential's ability to take large steps on
short horizons; however, the flux-adaptive method provides greater robustness on
longer horizons where the accept/reject mechanism degrades.

These results are presented in Tables~7--9 and Figures~7--8 of the revised
manuscript

\medskip

\noindent\textbf{Comment 2.6:} \textit{The authors should attempt to
numerically demonstrate the specific benefits of flux-based protection versus
pure probability pruning, and of $\Phi_{\max}$-based stepping versus other
activity indicators.}

\textbf{Response:} The comparison in the new Section~5.6 directly demonstrates
these benefits. The bottleneck system provides the clearest illustration of
flux-based protection: the flux-preserving method retains the $B = 1$ state
(probability $\sim 10^{-6}$, below any reasonable probability threshold) because
its outgoing flux is significant relative to the total system flux, producing
$\langle C \rangle = 357$ versus $\langle C \rangle = 18.7$ under
probability-only pruning. The Oregonator comparison demonstrates the benefit of
$\Phi_{\mathrm{total}}$-based stepping versus fixed or accept/reject stepping:
the flux-adaptive rule adjusts $\Delta t$ over two orders of magnitude
($10^{-6}$ to $10^{-4}$) to match instantaneous system activity, while the
accept/reject mechanism is unable to increase $\tau$ beyond $10^{-6}$ without
violating mass conservation, resulting in a 10$\times$ slowdown and incomplete
coverage. Quantitative metrics are reported in Tables~7--9

\medskip

\noindent\textbf{Comment 2.7:} \textit{Some guidance for choosing the key
tolerances $(\alpha, \varepsilon_{\mathrm{flux}}, \varepsilon_{\Delta t})$
across different types of models would be very valuable.}

\textbf{Response:}
% TODO: Parameter guidance table

%======================================================================
\section*{Summary}
%======================================================================

% TODO: Write summary once all responses are complete

\vspace{1em}
\noindent Sincerely,\\[0.5em]
Aditya Dendukuri, Shivkumar Chandrasekaran, and Linda Petzold

\end{document}
